\section{Demo: Camera Failures in Simulated Humanoid Dodgeball}
\label{app:demo}
\subsection{What viewers see}
The demo presents synchronized replays of the experiments in this paper.
Viewers compare camera conditions, pause at measured events, and follow the
robot's pelvis relative to the 1\,m boundary. The examples use the first declared
seeds; aggregate results are reported separately so that an illustrative trial
is not mistaken for average performance.

The overview video first compares healthy depth with camera freezes at ball
launch and 320\,ms later. In the illustrated trial, the later freeze allows a
dodge but the robot subsequently leaves the area. The video then compares
restoration of fresh images, followed by healthy, stale, and blank images for
$\pi_2$. The opening slide defines $\pi_1$ and $\pi_2$ and identifies their sources;
separate tables report the 64-seed fixed recovery test and the 192-seed input
comparison. These are two sets of learned weights trained with the same method
and recorded settings, not two proposed training methods.

The recovery replay compares continuously stale images with fresh
images restored at two or three seconds. All three robots dodge the same throw;
only the two-second restoration stays within the area in this example. The
three-second example is labeled exploratory. The
paper's aggregate recovery results, rather than this single replay, support
the reported effect.

\subsection{Playback and setup}
The videos use aligned views and synchronized holds so that the encounter and
later movement can be compared. An orange marker remains briefly after the
ball's native removal to make its last recorded position visible. The marker
is a display aid, not a simulated rebound or continued ball physics.

The demo consists of recorded simulation, rather than live policy execution.
A laptop with a browser is sufficient for presentation. The presenters bring
the laptop; a power outlet and an optional external display are useful.
Downloaded videos play offline. No physical robot, GPU, or special floor area
is needed. Online playback of the YouTube version requires internet access.

\subsection{Video}
The overview video is available at \url{https://youtu.be/MG4fqtBPVOY}.
Synchronized replays are available on the project
website linked below the paper title.

The results concern two fixed policies in one simulated G1 task. They do not
establish hardware safety or variation across independent training runs.

\subsection{Code, data and weight availability}
Evaluation code and configurations are available in the project's public
repository, linked from the project website. The four recorded-trial archives are prepared but have not been
publicly released. The released reference weights can be obtained from their
original source; the separately trained weights are not distributed.
Recomputing the reported statistics requires the trial archives. Rerunning all
simulations also requires the simulator assets and both sets of weights.
Public end-to-end reproduction is therefore not currently available.
